\chapter{Kommentierte Quellen}

Achtundzwanzig Werke, eingeordnet: Zu jedem steht, wofür man es aufschlägt und was es leistet, das andere nicht leisten.

\lit{1}{Bachelier, L.: \textit{Théorie de la spéculation}. Ann.\ Sci.\ École Norm.\ Sup.\ 17, 1900.}{Die Doktorarbeit, die die Brownsche Bewegung fünf Jahre vor Einstein beschrieb --- für Kurse statt für Teilchen. \textbf{Enthält bereits eine Optionsformel und das Spiegelungsprinzip.} Sie wurde ein halbes Jahrhundert übersehen; heute gilt sie als Gründungsurkunde der Finanzmathematik.}

\lit{2}{Bauer, H.: \textit{Wahrscheinlichkeitstheorie}. 5.\ Aufl., de Gruyter 2002.}{\textbf{Auf Deutsch und maßtheoretisch streng.} Bedingte Erwartung, Martingale, Konvergenzsätze mit vollständigen Beweisen. Die deutschsprachige Grundlage für alles ab Kapitel~7; knapp geschrieben und ohne Zugeständnisse.}

\lit{3}{Björk, T.: \textit{Arbitrage Theory in Continuous Time}. 4.\ Aufl., Oxford 2020.}{Die Preistheorie aus Kapitel~13, sauber aufgebaut --- Nachbildung, Maßwechsel, Hauptsatz, Zinsmodelle. Für Leser mit mathematischem Vorwissen die klarste Darstellung; für den Einstieg zu formal.}

\lit{4}{Brémaud, P.: \textit{Markov Chains}. 2.\ Aufl., Springer 2020.}{Ketten, Warteschlangen und Simulation in einem Band --- deckt die Kapitel~2 bis 5 vollständig ab, samt Metropolis-Verfahren und Kopplungsbeweisen. Viele durchgerechnete Beispiele.}

\lit{5}{Dayan, P., Abbott, L.: \textit{Theoretical Neuroscience}. MIT Press 2001.}{Die mathematische Neurowissenschaft, mit dem Modell aus Kapitel~14. \textbf{Zeigt, wie Poisson-Prozesse und Diffusionen zur Beschreibung von Spikefolgen benutzt werden} --- und wo sie zu einfach sind.}

\lit{6}{Doob, J.: \textit{Stochastic Processes}. Wiley 1953.}{Das Buch, das die Martingaltheorie begründete --- optionales Stoppen, Ungleichungen, Konvergenz, alles in der Fassung, die bis heute gilt. Schwer zu lesen und historisch unentbehrlich.}

\lit{7}{Durrett, R.: \textit{Probability: Theory and Examples}. 5.\ Aufl., Cambridge 2019.}{Das umfassendste Lehrbuch der Wahrscheinlichkeitstheorie mit Prozessen --- von Grenzwertsätzen über Martingale bis zur Brownschen Bewegung. Frei verfügbar; dicht geschrieben, mit vielen Übungen.}

\lit{8}{Ewens, W.: \textit{Mathematical Population Genetics}. 2.\ Aufl., Springer 2004.}{Die Diffusionstheorie der genetischen Drift aus Kapitel~14, vollständig --- Fixierungswahrscheinlichkeiten, Zeiten, Verweildauern. Die Fundstelle für alles, was dort nur angerissen wurde.}

\lit{9}{Feller, W.: \textit{An Introduction to Probability Theory and Its Applications}, Band I. 3.\ Aufl., Wiley 1968.}{\textbf{Das Buch, aus dem Kapitel~1 stammt.} Fellers Darstellung des Arkussinus-Gesetzes ist bis heute unübertroffen, ebenso seine Warnung vor der Anschauung. Enthält die Berechnung der Pólya-Konstanten. Kaum gealtert.}

\lit{10}{Gardiner, C.: \textit{Stochastic Methods}. 4.\ Aufl., Springer 2009.}{Das Handbuch der Anwender --- Fokker-Planck-Gleichungen, Näherungsverfahren, Rauschen in physikalischen und chemischen Systemen. \textbf{Behandelt den Unterschied zwischen Itō und Stratonowitsch aus Sicht der Naturwissenschaft} ausführlicher als jedes mathematische Werk.}

\lit{11}{Gerstner, W., Kistler, W., Naud, R., Paninski, L.: \textit{Neuronal Dynamics}. Cambridge 2014.}{Die moderne Fassung der Nervenzellmodelle aus Kapitel~14, mit Anpassung an Daten. Frei verfügbar. Zeigt, wie aus der Diffusion mit Schwelle ein brauchbares Modell echter Zellen wird.}

\lit{12}{Grimmett, G., Stirzaker, D.: \textit{Probability and Random Processes}. 4.\ Aufl., Oxford 2020.}{\textbf{Das ausgewogenste Lehrbuch zu diesem Band.} Ketten, Poisson-Prozesse, Warteschlangen, Martingale, Brownsche Bewegung --- alles mit Beweisen, aber ohne Maßtheorie zu verlangen. Mit einem eigenen Aufgabenband.}

\lit{13}{Hackenbroch, W., Thalmaier, A.: \textit{Stochastische Analysis}. Teubner 1994.}{\textbf{Auf Deutsch und vollständig.} Itō-Integral, Semimartingale, stochastische Differentialgleichungen --- die deutschsprachige Entsprechung zu [24], mit sorgfältigen Beweisen und historischen Anmerkungen.}

\lit{14}{Hull, J.: \textit{Options, Futures, and Other Derivatives}. 11.\ Aufl., Pearson 2022.}{Das Standardwerk der Praxis --- was mit den Formeln aus Kapitel~13 tatsächlich geschieht, welche Erweiterungen gebräuchlich sind und wo die Modelle versagen. Wenig Mathematik, viel Marktwissen.}

\lit{15}{Itō, K.: \textit{On Stochastic Differential Equations}. Mem.\ Amer.\ Math.\ Soc.\ 4, 1951.}{Die Arbeit, in der die Formel aus Kapitel~12 bewiesen wird. Kurz, klar und erstaunlich gut lesbar --- Itō schrieb ohne den technischen Apparat, den spätere Darstellungen aufbauen.}

\lit{16}{Karatzas, I., Shreve, S.: \textit{Brownian Motion and Stochastic Calculus}. 2.\ Aufl., Springer 1991.}{Das maßgebliche Werk zu Teil~IV --- Konstruktion, Pfadeigenschaften, Itō-Kalkül, Maßwechsel, lokale Zeit. Anspruchsvoll und vollständig; das Buch, das man aufschlägt, wenn man eine Aussage genau braucht.}

\lit{17}{Kimura, M.: \textit{The Neutral Theory of Molecular Evolution}. Cambridge 1983.}{Die Zusammenfassung der Theorie, die Kapitel~14 berührt. \textbf{Kimuras Argument ist im Kern das Ergebnis dieses Bandes}: Wenn $4Ns$ klein ist, entscheidet der Zufall, nicht die Selektion.}

\lit{18}{Kingman, J.: \textit{On the Genealogy of Large Populations}. J.\ Appl.\ Probab.\ 19, 1982.}{Der Koaleszent --- die Umkehrung der Blickrichtung von den Nachkommen zu den Ahnen. Ein kurzer Aufsatz, der die Auswertung genetischer Daten von Grund auf verändert hat.}

\lit{19}{Klenke, A.: \textit{Wahrscheinlichkeitstheorie}. 4.\ Aufl., Springer Spektrum 2020.}{\textbf{Auf Deutsch, umfassend, modern.} Von der Maßtheorie bis zur Brownschen Bewegung und darüber hinaus --- das deutschsprachige Gegenstück zu [7]. Als Nachschlagewerk für alle Grundlagen dieses Bandes geeignet.}

\lit{20}{Levin, D., Peres, Y., Wilmer, E.: \textit{Markov Chains and Mixing Times}. 2.\ Aufl., AMS 2017.}{Die Theorie der Mischzeiten aus Kapitel~3 --- Kopplung, Abschneideerscheinung, Kartenmischen. Frei verfügbar. \textbf{Das Buch, das erklärt, warum sieben Riffelmischungen genügen.}}

\lit{21}{Meintrup, D., Schäffler, S.: \textit{Stochastik --- Theorie und Anwendungen}. Springer 2005.}{\textbf{Auf Deutsch und anwendungsnah.} Enthält Markow-Ketten, Poisson-Prozesse und Warteschlangen mit durchgerechneten Beispielen aus Technik und Medizin --- die zugänglichste deutschsprachige Darstellung der Kapitel~2 bis 5.}

\lit{22}{Mörters, P., Peres, Y.: \textit{Brownian Motion}. Cambridge 2010.}{Die feine Struktur der Pfade --- Nullstellenmenge, lokale Zeit, Hausdorff-Dimension, Gesetz vom iterierten Logarithmus. Enthält den Satz von Donsker. Frei verfügbar; die Vertiefung zu Kapitel~10.}

\lit{23}{Norris, J.: \textit{Markov Chains}. Cambridge 1997.}{Die klarste Darstellung der Kettentheorie in diskreter und stetiger Zeit --- Klassifikation, Rekurrenz, stationäre Verteilungen, Ergodensatz, alles mit vollständigen Beweisen. Die Fundstelle für Satz 3.2.}

\lit{24}{Øksendal, B.: \textit{Stochastic Differential Equations}. 6.\ Aufl., Springer 2003.}{\textbf{Der freundlichste Zugang zum Itō-Kalkül.} Beweist alles Wesentliche der Kapitel~11 und 12 und wendet es auf Filterung, Optimalsteuerung und Bewertung an. Für Leser mit Analysiskenntnissen der beste Einstieg in Teil~IV.}

\lit{25}{Ross, S.: \textit{Introduction to Probability Models}. 12.\ Aufl., Academic Press 2019.}{Prozesse ohne Maßtheorie, mit sehr vielen Anwendungen --- Erneuerungstheorie, Warteschlangen, Zuverlässigkeit. \textbf{Das Buch für Leser, die rechnen wollen und auf Beweise verzichten können.}}

\lit{26}{Shreve, S.: \textit{Stochastic Calculus for Finance II}. Springer 2004.}{Die Finanzmathematik von Grund auf, aus einer Vorlesung entstanden. Sorgfältig, ausführlich, mit vielen Zwischenschritten --- die pädagogisch beste Darstellung dessen, was Kapitel~13 in einem Kapitel behandelt.}

\lit{27}{Steele, J.\,M.: \textit{Stochastic Calculus and Financial Applications}. Springer 2001.}{Ein Lehrbuch mit ungewöhnlich viel Gespür für die Frage, warum die Dinge so sind. \textbf{Die Erklärung, warum der Auswertungspunkt in Kapitel~11 alles entscheidet, ist hier am überzeugendsten} ausgeführt.}

\lit{28}{Williams, D.: \textit{Probability with Martingales}. Cambridge 1991.}{\textbf{Der schönste kurze Zugang zur Martingaltheorie.} Baut die Maßtheorie in achtzig Seiten auf und beweist dann Konvergenzsatz, optionales Stoppen und die starke Gesetzmäßigkeit --- alles über Martingale. Die Fundstelle für die bedingte Erwartung und die gleichgradige Integrierbarkeit.}

\begin{merksatz}
\boxtitel{Wenn Sie nur drei Bücher lesen}
\begin{center}
\renewcommand{\arraystretch}{1.6}
\begin{tabular}{@{}P{1.0cm}P{3.4cm}P{3.8cm}@{}}
\toprule
\textbf{[12]} & \textbf{Grimmett, Stirzaker} & \textbf{Für alles}\\
\textbf{[28]} & \textbf{Williams} & Für die Martingale\\
\textbf{[24]} & \textbf{Øksendal} & \textbf{Für den Itō-Kalkül}\\
\bottomrule
\end{tabular}
\end{center}

\begin{center}
\begin{minipage}{0.88\textwidth}
\kursiv{Grimmett und Stirzaker decken die Teile I bis III vollständig ab, ohne Maßtheorie zu verlangen. Williams zeigt auf zweihundert Seiten, wie viel aus dem Martingalbegriff folgt. \textbf{Und Øksendal führt in den stochastischen Kalkül ein, ohne den Leser mit Allgemeinheit zu erdrücken.}}
\end{minipage}
\end{center}

\begin{center}
\begin{minipage}{0.88\textwidth}
\kursiv{Nach Interessenlage: [9] Feller für die Anschauung, [20] Levin u.\,a.\ für Mischzeiten, [8] Ewens für die Populationsgenetik, [26] Shreve für die Finanzmathematik, [10] Gardiner für die Naturwissenschaften. Auf Deutsch: [19] Klenke und [2] Bauer für die Theorie, [21] Meintrup und Schäffler für die Anwendungen, [13] Hackenbroch und Thalmaier für den Kalkül.}
\end{minipage}
\end{center}
\end{merksatz}

\bigskip
\begin{flushright}
\itshape\color{grau}
Ein Verlauf zeigt Muster.\\
Die Verteilung zeigt sie nicht.\\
Erst beides zusammen sagt,\\
was der Zufall allein schon kann.
\end{flushright}
